\documentclass[11pt]{article}
\usepackage[a4paper,margin=25mm]{geometry}
\usepackage{amsmath,amssymb,amsthm,hyperref}
\emergencystretch=2em
\newtheorem{theorem}{Theorem}\newtheorem{lemma}{Lemma}
\title{Finite-dice profile fluctuations at every bounded endpoint scale}
\author{Complete all-scales sharpness theorem; independently reviewed}
\date{30 September 2026}
\begin{document}\maketitle

\begin{theorem}\label{main}
For every fixed $L>0$ there are positive constants $c_L,C_L$ such
that the following holds for all sufficiently large $m$.
For every scale $\lambda\in[m^{-1},L]$, there is a fully
permutation-fair family of $m+1$ finite uniform dice and a face of
one die whose comparison probabilities $p_1,\ldots,p_m$ satisfy
\[
 c_L\lambda\le s:=\sum_{i=1}^m(1-p_i)\le C_L\lambda,
 \qquad
 \max_i|m(1-p_i)-s|\ge c_L\sqrt{\lambda/m}.
\]
All face labels can be chosen distinct. No bound on their number is
asserted. The family may depend on both $m$ and $\lambda$.
\end{theorem}

The theorem matches the square-root term in
\path{natural_scale_endpoint_profile.tex} across all the displayed
scales, up to fixed multiplicative constants in the realized $s$.
It does not prescribe the exact numerical value $s=\lambda$.
The separate example in \path{../endpoint_profile_sharpness.tex}
has $s\asymp m^{-2}$ and error $\asymp m^{-1}$, forcing the
additive floor as well.

\begin{lemma}[A Gaussian node at the required scale]\label{node}
Let $Q=\lceil(m+1)/2\rceil$ and let $\sigma$ be the $Q$-node
Gauss--Legendre probability rule on $[0,1]$. For every
$\lambda\in[m^{-1},L]$ it has a node $a$ of weight $v$ with
\[
 c\lambda/m\le1-a\le C\lambda/m,
 \qquad v\ge c\sqrt\lambda\,m^{-3/2}.
\]
The constants can be absolute once $m$ is sufficiently large in
terms of $L$.
\end{lemma}
\begin{proof}
Number the nodes downwards from $1$, and choose
$j=\lceil\sqrt{m\lambda}\rceil$. Since $m\lambda\ge1$,
\[
 \sqrt{m\lambda}\le j\le2\sqrt{m\lambda}.
\]
For large $m$ depending on $L$, $j\le Q$ and $j/m\to0$
uniformly over the stated scale range. Write $2a-1=\cos\theta$.
The Legendre zero inequalities already used and justified in
\path{../endpoint_profile_fluctuation_barrier.tex} are
\[
 \frac{(j-1/2)\pi}{Q+1/2}\le\theta
       \le\frac{j\pi}{Q+1/2}.
\]
They imply $\theta\asymp j/m$, uniformly even for $j=1$.
Thus $1-a=(1-\cos\theta)/2\asymp j^2/m^2\asymp\lambda/m$.

For completeness, the Bernstein--Legendre bound from the same
reviewed note yields
\[
 \sum_{r=0}^{Q-1}(2r+1)P_r(2a-1)^2
     \le CQ/\sin\theta.
\]
The Gaussian probability weight is the reciprocal of this diagonal
reproducing kernel. Hence
\[
 v\ge c\sin\theta/Q\ge cj/m^2
      \ge c\sqrt\lambda\,m^{-3/2}.
\]
The two classical inequalities are recorded respectively in
\href{https://dlmf.nist.gov/18.16.E2}{DLMF 18.16.2} and
\href{https://dlmf.nist.gov/18.14.E7}{DLMF 18.14.7}; the cited
research note supplies the reproducing-kernel argument.
\end{proof}

\begin{proof}[Proof of Theorem~\ref{main}]
The Gaussian rule integrates every polynomial through degree $m$,
since $2Q-1\ge m$. Let $F$ be its CDF. Its two one-sided values
at the selected node $a$ differ by $v$, so at least one, denoted
$F_*$, satisfies $|F_*-a|\ge v/2$.
Choose $t\in(0,1)$ strictly between Gaussian nodes and sufficiently
close to $a$ on that side that
\[
 F(t)=F_*,\qquad |t-a|<v/8,\qquad
 |t-a|<(1-a)/10.
\]
Arbitrarily close one-sided approach makes all these requirements
compatible. Consequently
\[
 |F(t)-t|\ge3v/8,\qquad
 c\lambda/m\le1-t\le C\lambda/m.
\]

Consider the ideal row $(F(t),t,\ldots,t)$ and write
\[
 s_0=(m-1)(1-t)+(1-F(t)).
\]
Its lower bound is $s_0\ge c\lambda$. For the upper bound it is
enough to control the Gaussian tail; no cumulative quadrature-weight
asymptotic is needed. Exactness of the $m$th moment gives
\[
 1-F(t)=\sigma((t,1])
 \le t^{-m}\int_0^1x^m\,d\sigma(x)
 =\frac{t^{-m}}{m+1}\le\frac{C_L}{m}.
\]
Indeed $1-t\le CL/m$, so $t^{-m}\le\exp(C'L)$ for large $m$.
Since $\lambda\ge1/m$, this proves $s_0\le C_L\lambda$.
The distinguished coordinate's ideal deviation is
\[
 |m(1-F(t))-s_0|=(m-1)|t-F(t)|\ge c m v.
\]

Apply the reviewed realization theorem
\path{../finite_profile_realization.tex} to $\sigma,t$ with
\[
 0<\rho<\min\{v/100,\ c\lambda/(100m)\},
\]
where $c>0$ is small enough for the preceding lower bound on $s_0$.
Its hypotheses hold because $t$ is not a Gaussian node.
It produces an exactly permutation-fair family of finite uniform
dice with a comparison row satisfying
\[
 |p_1-F(t)|<\rho,\qquad |p_i-t|<\rho\quad(i\ge2).
\]
Thus $|s-s_0|\le m\rho$, preserving both comparability bounds.
Writing $e_i=p_i-t$ for $i\ge2$ gives
\[
 m(1-p_1)-s=(m-1)(t-p_1)+\sum_{i=2}^m e_i.
\]
Its absolute value is at least
\[
 (m-1)(3v/8-2\rho)\ge c m v
       \ge c\sqrt{\lambda/m},
\]
as required. Distinct labels and uniform face probabilities are
part of the realization theorem.
\end{proof}
\begin{theorem}[Sharp order at every bounded endpoint scale]\label{allscales}
Define
\[
 R_m(t)=\min\{mt,\sqrt{t/m}+m^{-1}\}\qquad(t\ge0).
\]
For every fixed $L>0$, the natural-scale upper theorem and the trivial
bound give, for every admissible monotone tensor and $s\le L$,
\[
 \max_i|m(1-p_i)-s|\le C_L R_m(s).
\]
Conversely, for every sufficiently large $m$ and every
$\lambda\in(0,L]$, actual finite uniform permutation-fair dice have a
comparison row with $c_L\lambda\le s\le C_L\lambda$ and
\[
 \max_i|m(1-p_i)-s|\ge c_L R_m(\lambda).
\]
As before, scale comparability is asserted, not the exact value
$s=\lambda$, and there is no bound on the face counts.
\end{theorem}
\begin{proof}
Since $0\le u_i\le s$, one has
$|m u_i-s|\le(m-1)s\le ms$ for $m\ge2$.
Combine this with the natural-scale upper theorem to prove the upper
bound. If $\lambda\ge1/m$, Theorem~\ref{main} supplies the lower
bound, since $R_m(\lambda)\le2\sqrt{\lambda/m}$.

For smaller scales use the Legendre construction from
\path{../endpoint_profile_sharpness.tex}, with adjustable amplitude.
Let $r\in\{m,m+1\}$ be odd and put
$\varepsilon_0=1/[8r(r+1)]$.
For $0<\varepsilon\le\varepsilon_0$, let $\sigma_\varepsilon$
have mass $\varepsilon$ at each endpoint and interior density
\[
 f_\varepsilon(t)=1-2\varepsilon P_r'(2t-1).
\]
The Legendre derivative bound gives $7/8\le f_\varepsilon\le9/8$,
and the interior CDF is
$F_\varepsilon(t)=t-\varepsilon P_r(2t-1)$.
Integration by parts and Legendre orthogonality show that
$\sigma_\varepsilon$ has the uniform moments through degree $m$,
exactly as in the cited note. Also
\begin{equation}\label{legendreclose}
 P_r(2t-1)\ge1-r(r+1)(1-t).
\end{equation}
The ideal comparison row $(F_\varepsilon(t),t,\ldots,t)$ has
\[
 s_0=m(1-t)+\varepsilon P_r(2t-1),\qquad
 |m(1-F_\varepsilon(t))-s_0|
       =(m-1)\varepsilon|P_r(2t-1)|.
\]

If $0<\lambda\le\varepsilon_0$, choose $\varepsilon=\lambda$
and $t=1-\varepsilon/(100m)$.
Equation \eqref{legendreclose} gives $P_r(2t-1)\ge1/2$, so
$s_0\asymp\lambda$ and the deviation is at least $cm\lambda$.
In this range $R_m(\lambda)=m\lambda$.

If $\varepsilon_0<\lambda<1/m$, choose
$\varepsilon=\varepsilon_0$ and $t=1-\lambda/(100m)$.
For all sufficiently large $m$, \eqref{legendreclose} again gives
$P_r(2t-1)\ge1/2$. Consequently
\[
 \lambda/100\le s_0\le\lambda/100+\varepsilon_0\le2\lambda,
 \qquad |m(1-F_\varepsilon(t))-s_0|\ge c/m.
\]
This is at least a fixed multiple of $R_m(\lambda)$, since
$R_m(\lambda)\le2/m$ throughout the present range.

In both cases $t$ is an interior continuity point of
$\sigma_\varepsilon$. Apply the finite realization theorem with
$\rho<\min\{\varepsilon/100,c\lambda/(100m)\}$, taking $c$
smaller than the fixed lower comparison constant for $s_0$.
The error in $s$ is at most $m\rho$, and the deviation of the
marked coordinate changes by at most $2m\rho$. Thus all preceding
lower and comparability bounds survive. The resulting dice are
exactly fair, equiprobable, and have distinct labels. This proves
the tiny-scale cases and completes the theorem.
\end{proof}
\paragraph{Sharp synchronization threshold.}
The natural-scale upper bound gives uniform relative synchronization
when $ms\to\infty$ on bounded endpoint ranges. This growth condition
cannot be replaced by a bounded-$ms$ assertion: choose
$\lambda=1/m$ above. Then $s\asymp1/m$ and the absolute deviation is
at least $c/m$, so $\max_i|m(1-p_i)/s-1|$ stays bounded away from zero.
\end{document}
